\chapter{Dashboard and Results}
\label{ch:dashboard}

\section{Dashboard Architecture}

The Streamlit dashboard is the only component the branch supervisor interacts with directly; every other service is invisible infrastructure. It must therefore surface the right information clearly while remaining responsive under continuous use.

\textbf{Why the dashboard cannot consume Kafka directly.} Streamlit re-runs the entire script from top to bottom on every user interaction --- clicking a button, changing a filter, navigating pages. A Kafka consumer embedded in the script would be destroyed and recreated on each interaction, losing its consumer group offset and replaying from the last committed position. Direct consumption is therefore architecturally incompatible with Streamlit. The solution is the sink consumer pattern of Chapter~\ref{ch:streaming}: two independent consumers write continuously to PostgreSQL and Redis, and the dashboard reads from PostgreSQL, which gives stable, queryable, filterable access to all historical alerts.

\textbf{Data access.} Reads go directly to PostgreSQL via \code{psycopg2} with Streamlit's \code{@st.cache\_data} decorator for query-level caching, with TTLs tuned per page: 10 seconds for the alerts list (freshness matters) and 30 seconds for KPI aggregations (stability matters more). A \code{safe\_query()} wrapper handles connection resilience --- if a query fails on a stale connection, common after Docker restarts or idle timeouts, it clears the cached connection, reconnects, and retries transparently. This single-retry pattern handles the most common failure mode without the complexity of a connection pool.

Writes go \emph{exclusively} through FastAPI endpoints; the dashboard never issues \code{INSERT} or \code{UPDATE} statements. This preserves the single-writer principle and keeps all write-path logic --- validation, audit logging, the feedback loop --- in one place.

\textbf{Refresh and navigation.} The \code{streamlit-autorefresh} component triggers re-execution every 10 seconds. Streamlit supports no native server-to-client push, so periodic polling is the pragmatic alternative; 10 seconds balances responsiveness (a new BLOCK appears within 10 seconds of being written) against unnecessary database load. The dashboard uses Streamlit's multi-page structure, so navigating executes only that page's script, and a shared \code{components.py} module renders the sidebar with pipeline status and pending alert counts by tier.

\section{Pages}

\textbf{Alerts} --- the supervisor's primary workspace. Alerts are filterable by tier, operation type, and review status, each row showing a color-coded tier badge, operation, amount, timestamp, fused score, and status. Expanding an alert reveals the full context for a decision: client, employee, and branch identifiers; triggered rules with descriptions; regulatory flags with warning icons; and the SHAP explanation.

The explanation is presented in \emph{two layers}. The first is the human-readable narrative from the deterministic NLG templates --- sentences such as \emph{``This versement is flagged because the amount is 4.2 times the client's 30-day average and it is the third deposit within 48 hours.''} The second, behind an expander, shows raw feature-level contributions: per-feature SHAP values for the Autoencoder and expected-versus-actual deviations for the LSTM. This serves both audiences --- supervisors needing a quick justification and analysts wanting the model's internal reasoning.

Each pending alert offers \textbf{Confirm} and \textbf{Reject} buttons with an optional comment. Either sends a POST to \code{/api/v1/feedback}, which updates the alert's supervisor decision and timestamp and upserts a record into \code{evaluation\_labels}. The page then clears its cache and re-renders immediately rather than waiting for the next autorefresh.

\textbf{KPIs} --- four metric cards (total alerts and the count for each actionable tier) above four Plotly charts: tier distribution as a pie chart, alerts grouped by operation type, alerts per day over the simulation period, and the top 20 branches by alert count. All aggregations use a 30-second cache TTL.

\textbf{Simulation} --- lets a supervisor construct a hypothetical transaction and score it without injecting it into the live pipeline. This serves two purposes: testing whether a suspicious transaction would alert before it is processed, and training supervisors to understand how the system evaluates different profiles. The form collects client, employee, amount, operation type, and conditional payload fields, validates them client-side, and POSTs to \code{/api/v1/score}, which runs the full pipeline synchronously --- enrichment, both models, fusion, SHAP, and decision rules --- returning tier, score, triggered rules, and the same two-layer explanation.

\textbf{Health} --- queries the Prometheus HTTP API for six service status indicators (three streaming services, two sinks, the simulation API) using the \code{up} metric, plus three operational metrics derived from the instruments of Chapter~\ref{ch:mlops}: throughput from the Counter's rate, p95 processing latency from the Histogram, and consumer lag per partition from the Kafka exporter. If Prometheus is unreachable the page degrades to a warning rather than crashing --- a monitoring outage should not block access to the other pages.

\section{Pipeline Validation and Results}

With all components integrated, the system was validated end-to-end by replaying the full 243{,}000-event dataset through the streaming pipeline and observing the resulting alert distribution.

\begin{table}[htbp]
\centering
\small
\caption{Alert tier distribution on the synthetic dataset}
\begin{tabular}{@{}llp{6.5cm}@{}}
\toprule
\textbf{Tier} & \textbf{Proportion} & \textbf{Interpretation} \\
\midrule
INFO & 89\% & Normal operations, no action required \\
REVIEW & 11\% & Mild deviations, logged for audit \\
ALERT & 0.35\% & Significant anomalies, supervisor review \\
BLOCK & $<$0.1\% & High-confidence anomalies, immediate action \\
\bottomrule
\end{tabular}
\end{table}

Scores span 0.015 (routine transactions by established clients) to 0.999 (anomalies matching multiple detection mechanisms). The distribution is heavily right-skewed, with the vast majority of events well below $T_1$ --- the expected shape for a functioning detector, where the tiers form a steep pyramid with very few events at the top.

\begin{table}[htbp]
\centering
\small
\caption{Model AUC-ROC on the 30-feature schema}
\begin{tabular}{@{}lcc@{}}
\toprule
\textbf{Model} & \textbf{AUC (74 features)} & \textbf{AUC (30 features)} \\
\midrule
Autoencoder & 0.9853 & 0.9714 \\
LSTM & 0.9866 & 0.9355 \\
\bottomrule
\end{tabular}
\end{table}

The decrease after feature reduction is expected and desirable. The 74-feature schema included raw profile statistics that made some anomalies trivially separable --- a client whose \code{tx\_count\_30d} is zero but who suddenly transacts is obvious to any model. The 30-feature contrast schema forces detection from \emph{relative} deviations rather than absolute values, which is harder but generalizes: the same features work across archetypes without archetype-specific thresholds. The LSTM loses more than the Autoencoder because it consumes sequences of 30-dimensional vectors, so each step carries less information per dimension, reducing its ability to model fine-grained temporal patterns. Both models nonetheless remain well above the 0.90 mark indicating reliable discrimination.

\textbf{Synthetic data caveat.} These results validate the system's ability to learn patterns, serve predictions at streaming latency, and produce a realistic alert distribution --- but they are not a production benchmark. The generator produces anomalies from known scenarios with known parameters, while real anomalies are messier, more diverse, and not drawn from the same distribution. The engineering argument is different from the statistical one: plugging in real data requires no code changes. The generator is replaced by real event sources, the batch pipeline rebuilds profiles from real history, the models are retrained on real features from the same \code{\_compute\_features()} function, and the retraining triggers ensure the models adapt as the distribution evolves. The AUCs confirm that this pipeline works end-to-end, not that it will reproduce these numbers on production data.
